\documentclass{article}
\usepackage[utf8]{inputenc}
\usepackage[T1]{fontenc}
\usepackage{geometry}
\usepackage{graphicx}
\usepackage{float}
\usepackage{hyperref}
\usepackage[english,indonesian]{babel}

\title{Optimasi Pemrosesan Template LaTeX Menggunakan Pendekatan Berbasis AI}
\author{Academic Support Agent}
\date{\today}

\begin{document}
\maketitle

\begin{abstract}

\end{abstract}

\chapter{PENDAHULUAN}


\chapter{TINJAUAN PUSTAKA}


% Original content (commented out):
% \documentclass[12pt]{article}

\usepackage[utf8]{inputenc}
\usepackage{latexsym,amsfonts,amssymb,amsthm,amsmath}

\setlength{\parindent}{0in}
\setlength{\oddsidemargin}{0in}
\setlength{\textwidth}{6.5in}
\setlength{\textheight}{8.8in}
\setlength{\topmargin}{0in}
\setlength{\headheight}{18pt}



\title{Math 333 Weekly Homework X}
\author{Your Name Here}

\begin{document}

\maketitle

\vspace{0.5in}



\subsection*{Exercise 1}
(Statement of problem goes here.)\\

\begin{proof}
(Type your proof here.)
\end{proof}

\vspace{2in} %Leave space for comments!


\subsection*{Exercise 2}
(Statement of problem goes here.)\\

\begin{proof}
(Type your proof here.)
\end{proof}

\vspace{2in} %Leave more space for comments!







\end{document}


\end{document}